\chapter{The Argument}
\label{chap:argument}

% Closed question: Which names must be spoken, and in what order, for the
% record to license public inference?

The thirty-seven names that follow are listed in the order the project's Lean
source has them. They are not philosophical chapters. Each is a class header,
a small public obligation that the next move in the argument has to satisfy
before it is allowed. The procession is one long inhale: a mark becomes
distinguishable, becomes admissible, becomes countable, becomes encoded, and
so on through residue, staging, overhead, transport, and decision. By the end
of the procession a public sentence may be inferred from the record. Before
the procession begins no such sentence may be spoken.

Two running examples carry the procession. Movements I through IV follow a
used-car inspection: a buyer, a seller, a mechanic, and a decision about
whether money should change hands. The inspection is ordinary. The names made
necessary by the inspection are not. At the close of Movement IV the example
promotes. The aircraft that the pre-flight checklist gated in Chapter
\ref{chap:scope} takes off, and Movements V and VI follow it through
transport, locality, halting, compilation, and the post-flight inference. The
car and the flight ask the same question at different scales.

A small numerical annotation rides in the right margin at section headers
where the reading has been verified against the current source. It is the
project's Bullshit meter: a counter the Lean source has been keeping all
along of how much overhead the argument has accumulated to carry itself. It
rises slowly through Movements I and II, accelerates through III, and
explodes through IV. No commentary on the number. The number is the gag, and
the gag is also the second argument of this chapter.

\section{\lean{Fact}}
\label{se:fact}
\meter{$\sim 9$}

Before anything can be distinguished, the record needs something it can treat
as the case. The very first thing a public process requires is a candidate
for entry. Without that, there is nothing to refine.

\lean{Fact} is the smallest public assertion. It is not yet measurement. It
is not yet evidence. It is the presence of a claim in a form the system can
carry. The system does not need to believe the claim. It needs to know the
claim exists and to be able to point at it later, by a name the system has
already agreed will refer to it.

The seller posts a 2017 sedan, 82{,}441 miles, four owners, automatic, clean
title claimed. The buyer has not driven the car. The mechanic has not opened
the hood. Nothing yet is admitted. But a claim has entered the record. The
listing exists, dated and tied to a phone number. From this point the
inspection has something to refuse, accept, or extend.

A receiving clerk stamps ``delivered'' on a package. The stamp does not
explain the contents and does not name the sender. It says a fact has
entered the office.

A fact that cannot be told apart from another fact cannot do work.

\section{\lean{DISTINGUISHABLE}}
\label{se:distinguishable}
\meter{$\sim 28$}

The first problem is not truth. It is difference. Before a ledger can say
what happened, it must be able to tell one mark from another. Before a ruler
can measure, it must have ticks. Before a court can hear testimony, it must
know which witness is speaking. Before any record can grow, the candidate
entry must be tellable apart from the entry that came before it.

\lean{DISTINGUISHABLE} is long because the obligation is long. It does not
say ``different.'' It says ``capable of being distinguished.'' The focus is
not on the object alone. The focus is on the object under a process that can
tell it apart from a peer. Two photographs of the same view, taken from the
same angle, with no chain of custody between them, are not distinguishable
from each other no matter how the optics differ. Two patients with the same
name and the same date of birth at the same clinic are not distinguishable
until a chart number, a wristband, or a barcode names them as separate. The
obligation puts pressure on the instrument, not on the world.

For the car, the VIN is the first instrument of distinguishability, but it is
not alone. The trim, the engine code, the title number, the key fobs, the
service records, the registration history, and the current owner together
pick out this car rather than a similar car from the same lot. Two 2017
sedans of the same trim and color are not distinguishable until the
instruments above have done their work. The mechanic photographs the VIN
plate against a date-stamped clipboard before the inspection begins.

Two unsigned, identical VIN-plate photographs are not distinguishable until
a chain of custody is named. Difference without witness is not difference
yet. The instrument refuses the entry.

A pharmacy distinguishes two patients with the same last name by date of
birth before filling a prescription. Without the date the prescription is
not yet ready to fill.

Not everything distinguishable should be admitted.

\section{\lean{ADMISSIBLE}}
\label{se:admissible}

Permission is not the same as difference. Noise can be distinguished. Errors
can be distinguished. Accidents can be distinguished. An instrument needs a
gate. A court needs rules of evidence. A ledger needs a standard for what
may enter and what may not. The gate is not the same instrument as the
distinguisher.

\lean{ADMISSIBLE} names the gate. It says the mark is not merely visible. It
is allowed into the record under the rules of the instrument. A reading that
violates the calibration window is admissible only if the calibration window
is widened, and the widening is itself recorded. A reading from a probe
whose calibration certificate has expired is not admissible. The gate's
strictness is part of the instrument's identity. Two instruments can produce
the same numbers under different admissibility rules and disagree about
whether the same physical reading entered the record or did not. The gate is
where the public commits to which kind of instrument it is dealing with.

A service invoice with shop name, date, VIN, and mileage is admissible. A
photograph of an undated parts box is not. A timing-belt service record from
a dealer's database is admissible. A seller's recollection that ``my cousin
said it was maintained'' is not the same kind of evidence. The mechanic's
inspection report admits the readings the mechanic took with the mechanic's
tools, under the mechanic's signature, on the mechanic's date. Everything
else is hearsay until it is brought to the same gate.

A handshake promise about prior maintenance is not admissible. No invoice,
no entry. The buyer does not have to be rude about it. The buyer has to be
correct about which kind of instrument the inspection is.

A tournament accepts scores submitted before the deadline on the official
card. Numbers texted after the round, even if accurate, do not enter the
record as scores.

Once admitted, the mark can be counted.

\section{\lean{COUNTABLE}}
\label{se:countable}

A mark that enters the ledger must be able to take a place. A pile of
admitted marks is not yet a record. The record has an order.

\lean{COUNTABLE} does not yet mean measured. It means the mark can be placed
into a sequence where ``next'' and ``previous'' make sense. The order may be
ordinal (this one came first), categorical (this one belongs in column B),
or ranked (this one matters more than that one). What is required is that
the sequence is defined, public, and stable. A countable record can be
pointed at by an index.

The inspection counts service-record entries in order: oil-change receipts,
tire replacements, brake work, timing-belt service, recall completion. The
count is not the diagnosis. The count is the first thing the buyer can ask
for. ``How many oil changes in the last fifty thousand miles?'' is a
countable question, and the answer determines what to ask next.

Numbered tickets at a bakery do not bake the bread. They make public order
possible. Customers can shop, walk away, and return; the bakery still knows
who is next.

A count needs a carried form.

\section{\lean{ENCODED}}
\label{se:encoded}

Counting has to survive transport. A count that exists only in one person's
memory is not yet a public count. The mark has to enter a form that someone
else can read.

\lean{ENCODED} names the entry of the count into a carried form. A count in
memory becomes a pulse in a wire, a tally on paper, a row in a database, a
mark in Lean. The encoding is not the count. It is the count's vehicle.
Different encodings make different things easy: a photo carries appearance,
a checkbox carries decision, a number carries magnitude. The encoding choice
already commits the system to what the next reader will be able to ask.

The car's condition enters forms: OBD-II trouble codes, inspection
checkboxes, tire-depth numbers, photographs of the underbody, service-record
PDFs, and a pricing spreadsheet. Each encoding is partial. Together they
travel from the shop bay to the buyer's email, to a lender's loan officer,
to a future mechanic if the car is later sold again.

A music score encodes rhythm and pitch so a performance can survive beyond
the composer's body. The score is not the music. The score is what lets the
music happen elsewhere and later.

Every encoding leaves a structured leftover.

\begin{movement}{From Entry to Residue}

The mark is now in admissible form. The procession has put a candidate into
the record, told it apart from neighbors, let it through the gate, given it
a place, and carried it into a form that can travel. The next pressure is
what survives the encoding without disappearing into it. Movement II begins
where every form leaves a structured leftover, and the instrument must
carry that leftover without pretending it is noise.

\end{movement}

\section{\lean{RESIDUE}}
\label{se:residue}

Not everything disappears into the form that carries it. The remainder
after division, the unread part of a signal, the unexplained error after
fitting, the obligation still open after a proof step: these are not
failures.

\lean{RESIDUE} names what remains and still matters. The encoding closed
neatly on most of the data and left a small piece outside. The piece is
small but not zero, and it does not behave like noise. It has shape. A
disciplined instrument records the residue separately, with its own label,
and refuses to absorb it into the headline reading.

The scan returns clean, but the steering wheel shakes under braking. The
title is clear, but one service interval is missing. The paint meter shows
a suspicious reading on one door. Each leftover is small. Each will become
an obligation that the inspection has to close, defer, or accept as a known
cost.

A bank reconciliation balances except for one outstanding check. The check
is not noise. It is structured residue, and the bookkeeper logs it as such.

Records must be compared.

\section{\lean{BINARY}}
\label{se:binary}

Comparison begins with a small question. Is this less than or equal to
that? Does this pass the threshold? Is this inside the permitted range?
The public interface is often a yes-or-no gate.

\lean{BINARY} names the comparison interface that produces a single bit.
The underlying physics may be continuous and the underlying record may be a
distribution, but the gate that fires at the end is two-valued. The bit is
the smallest signal the next stage can act on. The instrument is committing
to a decision the rest of the apparatus will use without re-doing the
comparison.

Brake pad thickness is above or below the shop's safety threshold. The
title is clean or branded. The asking price is within or outside the market
band the buyer agreed to before walking onto the lot. Each is a binary
that collapses an inspection or a search into a single answer the buyer's
next move can use.

A door sensor reports open or closed. The building automation system can
act because the state has a binary face. The door's hinge geometry is not
the sensor's concern.

A comparison that changes every time is not yet an instrument.

\section{\lean{REPEATABLE}}
\label{se:repeatable}

One good reading is not enough. If the same conditions do not yield the
same comparison, the instrument has not earned trust. Repeatability is not
truth. But without repeatability, truth has no place to stand. The first
reading that disagrees with the second is one of three things: a real
change, an instrument fault, or noise.

\lean{REPEATABLE} names the moral discipline of the instrument. A reading
that cannot be obtained again under the same conditions is hearsay, not
record. The instrument has to be able to be re-applied and the new outcome
has to match the old outcome within a public tolerance. The tolerance is
not optional. A claim of repeatability without a tolerance is a claim that
the instrument is perfect, which no instrument is. A claim of repeatability
without a stated condition is a claim that conditions do not matter, which
is also false. The instrument that earns the name has stated its conditions
and its tolerance and survived both.

The mechanic reads the battery under load twice, ten minutes apart, and
records both readings. The brake pads are measured at three points on each
pad. The road test is driven on the same route a second time after the
mechanic has had ten minutes to think. If a reading wanders without cause,
the buyer cannot use it; the mechanic logs the wander as a separate
observation and considers whether the instrument or the car is at fault.

A single high reading on a brake-pad gauge is not yet a measurement. One
observation is hearsay, not record. A mechanic who reports the high reading
as the brake-pad thickness has misrepresented what the instrument can do.

A tailor measures the same inseam twice before cutting cloth. The second
reading is not ceremony. It protects the cut, and the cut cannot be
undone.

A repeatable comparison can become a number.

\section{\lean{NUMERIC}}
\label{se:numeric}

A number is a carried comparison, not magic. The number does not arrive
from nature. It arrives from an instrument that has compared the world
against a reference, repeatedly enough to commit.

\lean{NUMERIC} names the moment when repeated comparison becomes a number
the record can use. The number is not the underlying quantity. It is what
the instrument has committed itself to as a public summary of the
comparison. Different instruments comparing the same world against
different references will produce different numbers, all of them honest.
The number belongs to the instrument as much as to the world.

The report says front brake pads are 4{\,}mm, battery health is 71
percent, compression is 178{\,}psi, and asking price is \$11{,}900. The
numbers are carried comparisons against thresholds, reference batteries,
factory specs, and market comparables. Each number has a unit and a
reference; without the reference the number does not yet mean anything to
the buyer.

A speedometer's displayed value is not speed itself. It is a numeric
output produced by a calibrated comparison of pulse count and time. The
driver trusts the comparison without re-doing it.

A number must stand for something unresolved or real.

\section{\lean{REPRESENTABLE}}
\label{se:representable}

The world contains things not yet closed. A lab result is pending. A
contract leaves a blank. A compiler carries a metavariable. A detector
leaves a gap between events. The unresolved thing still has to be carried.

\lean{REPRESENTABLE} says the unresolved thing can be carried in the
record as itself, without being prematurely resolved. The blank stays a
blank, but it has a name, a type, and a position. Later work can fill it
without re-deriving where it should sit. A representation that requires
resolution before it can be carried is a worse instrument than a
representation that can carry the unresolved item without spending it.

The buyer does not know whether the vibration is warped rotors, worn
bushings, tire imbalance, or frame damage. Each possible cause has a name
on the report. The shop schedules a follow-up if the buyer wants to
purchase. The vibration is carried as a represented obligation: tested by
one of four named diagnostic paths, none yet executed.

An architect marks ``future stair'' on a renovation plan. The stair is not
built, but it has a place, a type, and a clearance the design can carry
without committing to the construction.

Representation cannot float free as mere notation.

\section{\lean{PHYSICAL}}
\label{se:physical}

Symbols need contact with an instrument. A representation that lives only
as text is not yet a physical record. The instrument has to be able to be
tied to a process that could enter a measurement.

\lean{PHYSICAL} says the representation is anchored to a process whose
records the instrument can in principle examine. The anchor is what
separates speculation from claim. ``Frame damage'' as a sentence on a
report is not yet physical. ``Frame damage measured at lift points three
and five, paint depth at door two, weld signature at the rear subframe''
is physical: the sentence names processes that could be re-executed.

The suspicion of frame damage cannot remain a suspicion on the inspection.
It must touch the lift, the gap gauge, the paint meter, and the alignment
machine. The shop's report names each of these processes against each of
the suspect locations, and the result becomes a physical record the buyer
can take to a second shop if the buyer wishes.

A weather icon means little unless it is tied to barometric pressure,
radar return, or thermometer readings. The icon is the report; the
readings are what make the report physical.

Physical records must be brought under shared questions.

\section{\lean{COMPARABLE}}
\label{se:comparable}

Two records do not compare themselves. Different instruments, different
contexts, different units, different representations: comparison requires
a shared question and a shared frame in which the question can be asked.

\lean{COMPARABLE} names the permission to bring records under the same
comparison. The two records have been re-expressed in a common frame, or
the question has been reformulated so the frame difference does not
matter, or both. Without the permission the comparison is a category
error wearing the costume of arithmetic. The work to earn the permission
is part of the record, not preamble to it.

A dealer price, a private-party listing, an auction value, and an
independent repair estimate become comparable only under the buyer's
question: is this car worth this price after likely repairs? The four
numbers were collected under four different conditions, and a comparison
that ignored that would mislead. The buyer's question normalizes them.

Two recipes become comparable when both are converted to servings,
grams, oven temperature, and time. Before that conversion the two
recipes are disagreement in different units about different dishes.

Some comparisons cross into observation.

\section{\lean{OBSERVED}}
\label{se:observed}

Possible is not yet recorded. An event that could have been observed has
not yet been observed until the instrument has captured it and the
capture has entered the ledger. The line between possible and observed is
a public threshold the instrument has to cross.

\lean{OBSERVED} names the boundary crossing. The event has entered the
ledger. The instrument has admitted it. The record now contains the event
under a name, with a time, under a witness. Possible events are
infinite; observed events are not. The instrument's discipline is in
which possible events it allows itself to call observed, and in how it
logs the ones it refuses.

During the test drive the mechanic places a stethoscope on the intake
manifold and a ticking enters the report. The concern crosses from
possible sound to observed entry. The ticking is now in the record at a
named RPM, in third gear, on the test loop, at a recorded moment. The
buyer can ask the next shop to verify the same observation under the
same conditions.

A teacher does not grade a student's likely method. The answer becomes
observed when it appears on the page, and only then can be marked right
or wrong.

Some recorded things remain unresolved but present.

\section{\lean{PRESENT}}
\label{se:present}

Unresolved does not mean absent. A sealed envelope is present without
being read. A legal blank is present without being filled. A compiler
obligation is present while still unresolved. The record has to be able
to hold these without forcing closure.

\lean{PRESENT} refuses the lazy alternatives. The thing is not gone. It
is held open under a name. The instrument carries it forward to
subsequent stages without pretending it has been resolved and without
losing track of it. Things that are present but unresolved are different
from things that have been resolved and from things that never entered
the record; the discipline keeps the three categories separate.

A pending misfire code is present even if the check-engine light is off.
The buyer should not call the engine failed, but the buyer cannot pretend
the code is absent. The report logs the code as present, dated, under the
scan tool's serial number, with the relevant freeze-frame data, and notes
that no condition has yet been concluded from the presence.

An unsigned contract clause is present on the page. It is unsettled, not
nonexistent, and the contract cannot be filed as complete until the
clause is either signed, struck out, or formally deferred.

The open place must still be measurable.

\section{\lean{MEASURABLE}}
\label{se:measurable}

Holding open is not the same as refusing discipline. A pending obligation
that no test can address is not yet a useful obligation. The instrument
has to commit to a procedure that could resolve it.

\lean{MEASURABLE} says the unresolved place can still be asked
disciplined questions. It can be bounded, tested, partially resolved, or
proven unresolvable by current instruments. The instrument names the
questions it would answer if it had the time and the equipment. The open
place is not an invitation to speculate forever; it is a queue of tests
that have been designed but not yet executed.

The pending misfire can be tested by a compression test on the suspect
cylinder, a spark test, a fuel-trim review from the OBD data, a coil
swap to see if the code migrates, and a cold-start drive the next
morning. The buyer now knows what the next visit costs and what each
step would buy in information.

A cracked ceramic bowl can be tapped, lit, magnified, and weighed before
a conservator decides whether it can be repaired. None of those tests is
the repair; each is a measurable next step that constrains the decision.

Unresolved states need a process that carries them without pretending
they are resolved.

\begin{movement}{From Residue to Staging}

The leftover has been held open, named, and brought under tests the
instrument can name in advance. The record now contains both the closed
entries and the open obligations, each under its own discipline. The
next pressure is how a serious system moves such material through
stages. The held-open state has to travel from inspection bay to repair
queue to purchase decision without losing its provenance.

\end{movement}

\section{\lean{GUNGAN}}
\label{se:gungan}

Every serious system needs a way to carry the awkward thing. The misfire
is not failure. The shake on braking is not safety violation. The
missing service receipt is not fraud. They are all in the awkward
middle, between acceptable and unacceptable, and the system has to keep
them visible without forcing premature judgment.

\lean{GUNGAN} is a rude name for a real function. It marks the fact that
unresolved states still need social machinery, code machinery, and
rhetorical machinery to keep them from collapsing into either nonsense
or premature certainty. The polished name for this function would be
something like ``pending-state carrier'' or ``conditional admission
record.'' The polished names would let the function hide. The system
kept the rude name because the function is awkward, the cases that
require it are usually embarrassing, and the carriers themselves often
complain about being asked to hold a state instead of being given an
answer. The name calls the function by what it actually is.

The inspection report needs a status for ``hold for second opinion,''
``requires lift inspection,'' ``seller must provide receipt,'' and
``buyer may rescind if X is not provided by date Y.'' The shop's
software has checkboxes for each. The checkboxes do not have names like
``pending condition.'' They have names that match what the mechanic says
out loud: ``ask the seller,'' ``come back Thursday,'' ``do not buy
unless.'' Each checkbox is a small \lean{GUNGAN}: a held state with a
path out, but not yet a closure.

A library keeps damaged books on a repair cart instead of returning them
to circulation or throwing them away. The cart is not glamorous and it
embarrasses the catalogers. It is also the only thing standing between
the library and either a permanent loss or an unsafe loan.

The carried unresolved state can now pass through stages.

\section{\lean{SOURCE}}
\label{se:source}

A plan is not the act. A recipe is not the cake. A score is not the
music. A source file is not the running program. The system has to
distinguish what is staged from what has been performed.

\lean{SOURCE} names the form before performance. The source carries
sufficient detail that a competent performer could execute it, but the
source has not been executed yet. The source can be reviewed, signed
off, amended, or refused before any action is taken. Whatever the system
later performs has to be derivable from a source it can point at;
otherwise the performance cannot be reviewed against intention.

The mechanic's checklist is source as document. It names VIN, title,
scan, fluids, tires, brakes, frame, test drive, pricing, and final
recommendation while still sitting on paper or screen. The buyer can
review the checklist before authorizing the inspection. The checklist
names what would be done and refused; it is not yet the inspection
itself.

A theater script is source. It can name entrances, lights, and lines
without yet producing the performance. The performance derives from the
script and remains comparable to it.

The source must be put through a process.

\section{\lean{EXECUTED}}
\label{se:executed}

At some point the plan spends itself in action. Source becomes
performance. The instrument's source had said what would be done; now
the instrument has done it.

\lean{EXECUTED} says the record has crossed from source into
performance. Something has happened under the named process. The
execution itself is recorded, with its time and its operator, and the
recorded execution is different from the source: the execution can be
reviewed against the source as fidelity, against the world as outcome,
and against the buyer's expectation as adequacy.

The mechanic's body and tools meet the car: the lift arms settle under
the frame, the scanner plugs in, the brake pedal is pressed, the road
test is driven. The checklist changes medium from document to action.
Each item on the source list now has an executed counterpart, with a
time, a tool reading, and the mechanic's initials in the corresponding
box.

A recipe is executed when ingredients are weighed, mixed, heated,
cooled, and plated. The recipe still exists on the card; the dish is
the executed counterpart.

Execution must leave a carried result.

\section{\lean{VALUE}}
\label{se:value}

After execution, something must remain usable. The point of running the
instrument was not the running; it was the result the instrument
produced and the system can carry forward.

\lean{VALUE} is not vague importance. It is what the process can carry
forward after the source has passed through the machine. The value is
typed, named, dated, and distinguishable from the source that produced
it. A process whose only outcome was the execution itself, with no
carried result, is not yet doing the work the instrument was built for.
The value is the public deliverable.

The carried result is not ``the mechanic looked at it.'' It is the brake
reading at 4{\,}mm, the code list, the leak note, the repair estimate
at \$1{,}800, and the buy/walk/renegotiate recommendation. Each is
typed. The buyer can show any of them to a lender, a second shop, or a
friend without re-running the inspection.

After a blood-pressure cuff runs, the value is the reading, not the
nurse's effort. The chart records the reading; the nurse's effort is
implied by the chart's existence.

Value invites the question of how much.

\section{\lean{MAGNITUDE}}
\label{se:magnitude}

Quantity arrives late. The public mistake is to think measurement begins
with how much. It does not. The system has already had to distinguish,
admit, count, encode, compare, represent, observe, hold open, stage,
execute, and carry value.

Only then can it ask for \lean{MAGNITUDE}. The number arrives in a
frame the procession has built: a number against a named unit, against
a named reference, after a named execution, by a named instrument.
Asked too early, the magnitude is grandstanding. Asked at the right
place in the procession, the magnitude is the moment when the
instrument starts paying out its promised information.

Magnitude is the arrival of numbers at all: 4{\,}mm of brake pad,
178{\,}psi compression, \$1{,}800 of estimated repairs, 82{,}441 miles,
three drops of oil leak per week. Each number stands in the frame the
prior names have built. The buyer can compare numbers to other numbers
because the prior names made them comparable.

A kitchen scale gives 312 grams of flour. The magnitude arrives after
the bowl, the tare, the scale, and the ingredient have been
distinguished, admitted, and weighed in that order.

Magnitude without a scale is not enough.

\section{\lean{SCALED}}
\label{se:scaled}

How much, relative to what? A magnitude without a scale is just a
number. The scale is the public statement of which threshold turns the
number into a judgment.

\lean{SCALED} says the magnitude is not naked. It carries a unit, a
basis, a ruler, or a calibrated threshold. The same magnitude under
different scales yields different judgments, and the scale's identity
is recorded. Two instruments that report the same number but used
different scales did not report the same reading. The scaling is part
of the answer; the answer without the scaling is incomplete.

4{\,}mm becomes ``near service limit'' only against the brake-pad
threshold the shop's safety standard specifies. \$11{,}900 becomes
``too high'' only against comparable sales, against repair costs, and
against the buyer's budget. The mechanic states the threshold and the
comparables explicitly; the report would not be usable without them.

38.5 degrees Celsius becomes ``fever'' because the clinical threshold
is 38.0. The thermometer's reading is the same regardless of the
scale; the diagnosis depends on the threshold the clinic has adopted.

Carrying a scale creates load.

\begin{movement}{From Staging to Overhead}

The record now has scale. The instrument has carried a candidate from
raw listing through admissibility, residue, staging, and value, and
produced numbers a buyer can act on. The cost of all that carrying has
been the quiet subject of the meter in the margin and is about to
become the loud subject of the next movement. The leftover the system
has been accumulating has a name. The next movement opens its accounts.

\end{movement}

\section{\lean{LOAD}}
\label{se:load}

A scale is a burden the system must carry. To measure under a scale,
the system has to keep the scale available: the reference, the
calibration, the trained operator, the maintenance schedule, the
documentation.

\lean{LOAD} names that burden. The instrument does not get the scale
for free. Someone has paid for the buffer, the certificate, the
training, the storage, and the time. Some of the load is visible at
the moment of measurement; some sits backstage, in costs the
institution incurred to be the kind of place that can measure this
thing at all.

The buyer carries repair manuals, price guides, the cost of the
pre-purchase inspection itself, lender terms, state title rules,
registration fees, insurance quotes, and the mechanic's labor rate.
The scale costs attention, money, and a Saturday morning. None of
those costs is the car; all of them are required to buy the car
responsibly.

A laboratory's pH readings depend on buffers, calibration logs, clean
probes, and trained hands. The pH meter itself is the cheapest part of
the instrument's actual load.

Finite approximation may be enough.

\section{\lean{FINITE\_ELEPHANT}}
\label{se:finite-elephant}

Enough can be large and still finite. The buyer cannot inspect every
bolt, relay, gasket, and prior decision. The list of things that could
be checked is enormous. The list of things that will be checked has to
be bounded.

\lean{FINITE_ELEPHANT} is funny because the elephant is big. It is
serious because the elephant is finite. The name says a finite
approximation can carry enough of the structure for the next move to be
admissible. The list of checks is not exhaustive; it is sufficient
under a stated standard. The sufficiency is itself recorded, and the
standard is reviewable.

The buyer commissions the shop's standard pre-purchase inspection. The
checklist has roughly 150 items, not 15{,}000. The 150 items are large
enough to support a purchase inference under a written standard, and
small enough to be completed in three hours. The remaining 14{,}850
possible checks are not refused; they are out of scope, and the scope
is named.

A museum cannot examine every pigment molecule in a painting, but a
finite conservation survey can still support a treatment plan. The
survey's scope is the discipline that makes the plan reviewable.

Once enough is measured, overhead becomes visible.

\section{\lean{BULLSHIT}}
\label{se:bullshit}
\meter{$2733$}

Overhead is real. Paperwork, ceremony, friction, handwaving,
credentialing, performance, wasted motion, social pressure, repeated
explanation: the world is full of cost that is not the object but
attaches to the object. The instrument has to be able to name this
cost without becoming cynical.

The code names it \lean{BULLSHIT}. The name is funny because it is
impolite. It is accurate because it measures overhead. The class is
declared with a comment that says the quantity is strictly conserved:
it does not vanish from a system; it migrates. The polished
alternative would be ``accumulated operational overhead index.'' The
polished name would let the cost slip behind a clinical phrase, which
would let the cost grow quietly. The rude name keeps the cost visible.
The meter in the right margin of this section is the running total.
Every section the procession has crossed has added to it. The reader
has watched the number climb without commentary. The number is the
gag, and the gag is also the chapter's quietest argument.

The seller's ``never abused,'' the dealer's prep fee, the glossy
detailing, the warranty brochure, the urgency script (``I have two
other buyers this weekend''), and the financing theater do not fix the
car. They still affect the cost of reaching a decision. The buyer is
being asked to pay for them, either in money or in attention. The
inspection report carries an ``optional add-on'' line with the
dealer-suggested upcharges. The line is honest only when each upcharge
is labeled with what it actually buys, and many of them buy attention
diverted from the inspection.

A meeting about the meeting can consume the budget before the project
begins. The meeting is not the project. The cost of holding the
meeting is nevertheless deducted from the project's budget, and the
project's deliverable is whatever the budget can still afford after
the meetings have taken their share.

Overhead can reproduce itself as belief.

\section{\lean{PROPAGANDA}}
\label{se:propaganda}

Some overhead wants to convert the record around it. The phrase that
began as marketing becomes the phrase the buyer hears from friends,
then from neighbors, then from a forum thread that quotes another
forum thread that quotes a magazine review that quoted the original
press release. The record is not the source of the conviction.

\lean{PROPAGANDA} names the process by which a claim begins carrying
social force instead of merely evidential force. The claim may be
true. The problem is not truth-value. The problem is that the
conviction the public holds about the claim is no longer proportional
to the evidence the record supports. Conviction has been amplified by
repetition, by authority, by aesthetic packaging, or by social
pressure. The instrument that respects the difference between
evidential and social force can still act on the claim. The instrument
that confuses the two will buy the wrong car, vote against the wrong
proposition, or refuse the right treatment. The discipline is to keep
evidence and conviction on separate ledgers.

``One owner,'' ``highway miles,'' ``mechanic owned,'' ``rare package,''
and ``the last one of these to ever come up'' circulate as sales force
long after the records stop supporting them. The buyer hears the
phrases on the forum before walking onto the lot. The seller does not
even have to repeat the phrases; the buyer arrives already converted.
The inspection's job is to put each phrase back next to the records it
actually relies on and to show what the records say.

A phrase repeated often enough has not become evidence. The listing's
confidence is refused when the records do not carry it. The buyer's
confidence is also refused, by the same gate, when the records do not
support it either.

A product package saying ``doctor recommended'' may convert authority
into purchase pressure without showing the study. The packaging is the
propaganda; the study is the evidence.

Reproduced belief needs initiates.

\section{\lean{ACOLYTE}}
\label{se:acolyte}
\meter{$7699$}

A reproduced system recruits carriers. The slogan needs people to say
it. The people who say it the most fluently are the ones for whom
saying it has become identity. They are not lying. They are repeating,
with sincerity, claims that the record under their feet does not
support.

\lean{ACOLYTE} is the person or process initiated into the structure.
The name is not polite. It should not be. It marks the point where
overhead has become a social machine: a community has formed around
the claim, the community rewards fluency with the claim, and the
community's members acquire status by repeating the claim
convincingly. The acolyte is not necessarily a villain. The acolyte is
often a competent professional who has been trained to trust a phrase
the way an apprentice is trained to trust a teacher's instinct. The
problem is that the phrase is not the teacher's instinct. The phrase
is what the community pays its members to repeat, and the payment
makes the repetition feel like understanding.

The buyer may start repeating forum slogans: ``these engines run
forever,'' ``never buy the first model year,'' ``dealer service means
perfect,'' ``manual transmission means an enthusiast owner.'' The
slogans may contain experience. They may even contain truth in some
cases. They are not the inspected record of this car. The buyer who
walks into the inspection already saying these things has to decide
which sentences will survive contact with the record and which will
not. The mechanic's job is gentler than refusing the slogans: it is
asking, for each slogan, whether the record under this hood agrees.

A junior analyst repeats a firm's house phrase about a market because
the phrase is rewarded internally, not because the model has been
checked. The analyst is not dishonest. The analyst is doing what the
firm pays for. The firm's clients are still owed the model.

After propaganda and its acolytes, science must earn its name.

\begin{movement}{From Overhead to Truth Transport}

The procession has crossed the overhead movement. Cost has been named,
the finite elephant has been bounded, the rude names have done their
rude work, and the social machinery around the record has been put
back in its place. The used-car example has carried the procession
from entry to overhead. It cannot carry the rest. The names that
follow ask for transport across distance, contact with the real,
locality, universality, halting, and compilation; they ask for an
instrument that travels. The airplane that the pre-flight checklist
gated in Chapter \ref{chap:scope} has finished its own pre-flight.
The flight takes off here.

\end{movement}

\section{\lean{SCIENTIFIC}}
\label{se:scientific}
\meter{$23964$}

The word ``scientific'' cannot mean ``our side won.'' It cannot mean a
person in a lab coat said it. It cannot mean the story became popular.
It must name a process that lets the record learn while remaining
answerable to what was observed.

\lean{SCIENTIFIC} has to earn itself after \lean{PROPAGANDA}. The name
points to a discipline of recording, replicating, comparing, and
revising that does not depend on the social status of the recorder.
The discipline is not glamorous and it is often slow. It is the only
thing standing between an honest mistake and a permanent one.

The flight earns the word scientific through procedures that do not
depend on the captain's personality: clearance recorded by ATC, flight
plan filed and acknowledged, weather briefing logged, fuel computation
signed, weight-and-balance posted in the cockpit, departure stripped
from the tower list. None of the steps trusts the crew's mood; each
step has a record the next shift can read.

A seed trial is scientific when plots, controls, weather, soil, yield,
and replication are recorded so that the field can answer back to the
hypothesis.

Learning raises the question of truth.

\section{\lean{TRUTH}}
\label{se:truth}

Truth cannot float free. The instrument has carried readings,
comparisons, representations, observations, scales, and overhead. At
some point a public sentence has to be allowed to say: this is the
case. The sentence is dangerous if it forgets what carried it.

In this project, \lean{TRUTH} is truth carried by a process. The
system does not get to say ``true'' without saying how the truth
moved. The sentence is licensed by the carrying. The same words spoken
without the carrying have a different content; they are claims about
belief or authority rather than claims about record. A community that
says ``true'' without naming its carrying has slipped into propaganda
even when the underlying claim happens to be correct. The discipline
is to keep the sentence and its carriage joined, and to let the
joined object be the public artifact. The reader who hears the
sentence may then ask, at any time, what the carriage was. The
sentence is honest because the carriage is inspectable.

``The flight may be cleared at flight level three-five-zero'' is not a
naked truth. It carries the active weather, the en-route winds, the
fuel reserves, the aircraft's certified ceiling, the operator's
authorizations, the route's available altitudes, and the controller's
coordination with adjacent sectors. Each is named on a strip, in a
system, with a timestamp. The clearance is true under that carriage.
The same sentence in a simulator with the same words is not the same
truth: the carriage is different.

``The bridge is safe'' is truthful only through inspection date, load
rating, material condition, and engineering standard. The same
sentence on a sign without those carriers is decoration.

Movement of truth requires witness.

\section{\lean{WITNESSED}}
\label{se:witnessed}

A claim should not merely appear. The system that lets a true sentence
be spoken has to be able to say who carried the sentence, from where,
under what constraint, by what authority. Without witness, the
sentence has no return address; it cannot be challenged, verified, or
revised.

\lean{WITNESSED} says the claim passed through a recordable transport.
It was carried from somewhere, by something, under a constraint. The
constraint is what separates witness from rumor. A witness is not just
a person; it is a person operating under a public discipline, whose
testimony enters a record the discipline maintains. The witness's
authority is bounded by the discipline. A witness who steps outside
the discipline has stepped out of the role; the testimony from that
step is worth what hearsay is worth. The discipline that lets the
witness's testimony count is the same discipline that lets it be
challenged and, where necessary, refused.

The flight is witnessed by the ATC recording, the transponder, the
ACARS data link, the flight data recorder, and the cockpit voice
recorder. Each records a different aspect of the flight under its own
discipline. None of them is the flight. Together they constitute the
witness record that lets a post-flight inquiry, a maintenance review,
or a regulatory audit proceed without having to interview the crew.
The witness records outlive the flight by years.

A signature without a date is not witnessed enough for this transport.
The claim has a name, but no time. A clearance read back over a radio
without identification is also refused: the controller cannot file
what the controller cannot attribute.

A notary's seal does not make a statement true by magic. It witnesses
the identity and act under a public rule. The seal's power is the
rule, not the wax.

Witnessed truth must meet a world.

\section{\lean{REAL}}
\label{se:real}

The carried truth must answer to contact. A sentence that is
internally consistent but does not meet the world is decoration. The
instrument has to be able to point at the moment when the truth was
tested against something outside the record.

\lean{REAL} names the pressure that the claim meet the world rather
than only the syntax that carried it. The aircraft is or is not on
the runway. The patient is or is not breathing. The bridge does or
does not carry the load. The world's answer is not always immediate,
but it is unavoidable. The instrument's discipline is to keep the
test moment visible in the record.

The aircraft starts cold or it does not. It accelerates to V1 in the
expected distance or it does not. The wheels touch the runway at the
expected point or they do not. Paper meets metal at each of these
moments. The flight data recorder logs the metal's answer to the
paper's prediction; the post-flight review can compare them.

A restaurant reservation is real when the table exists, the time is
held, and the diners can be seated. A confirmation email without a
table at the hour is reservation theater.

Contact happens somewhere.

\section{\lean{LOCAL}}
\label{se:local}

Every reading has a place. A ruler sits somewhere. A clock ticks
somewhere. A compiler elaborates in a specific environment. A witness
sees from a position. The instrument has to be honest about where it
was when it read.

\lean{LOCAL} names honest placement. The Lean code carries a sidecar
here: \lean{BigRedDogProcess}. It is the project's local-derivative
process, parked inside \lean{LOCAL} rather than given its own class.
The name is rude on purpose; keeping a derivative attached to its
neighborhood has the cartoonish clumsiness of a very large, very
specific dog who can only stand in one backyard. The joke says with
less dignity what \lean{LOCAL} says with more.

The cockpit altimeter reading is local. It depends on the local
barometric setting, which is itself local to the airport's elevation
and the day's pressure. The same altimeter, flown the next day or in
the next country, reads against a different local reference. The
controller cannot accept ``altitude'' without its place. The altitude
clearance is given against a named reference at a named station.

A soil test from one field corner cannot speak for the whole farm
without a sampling plan. The test is correct about the corner; the
corner is not the farm.

Local readings want transport.

\section{\lean{UNIVERSAL}}
\label{se:universal}

People love to say ``always'' too early. The first reading from a new
instrument tempts the operator into universal claims. The discipline
is to hold the claim to what the instrument and the procession have
actually earned.

\lean{UNIVERSAL} is not the erasure of locality. It is the ambition to
transport local readings without losing admissibility. The universal
claim is what survives when the local conditions have been varied and
the reading still holds, or when the local conditions have been
measured and the variation has been accounted for. The universal
sentence keeps the local sentences underneath it; the universal does
not absorb the local.

The flight obeys airspace rules that apply to all aircraft of its
class in that airspace. The rules are universal in scope. The flight,
however, is one specific aircraft on one specific date, complying
with the universal rules at one specific point. The universal claim
about airspace is correct; the airspace remains itself; the specific
flight is the specific flight.

A building code transports lessons across buildings, but each
inspection still happens at one address. The code's universality is
not a substitute for the local inspection.

Universal comparison requires logic with calibration.

\begin{movement}{From Truth Transport to Decision}

The truth-bearing record has been carried with its witness and place.
The flight has crossed altitudes, talked to towers, made contact with
the real, and stayed local enough to be honest while universal enough
to be recognized. The procession is now at the edge of action. The
next pressure is what structure lets the process stop, cross an
instrument, survive translation, and finally license an inference
the public can hear.

\end{movement}

\section{\lean{LOGICAL}}
\label{se:logical}

Cleverness is not logic. A sharp argument with no scale is rhetoric.
The final movement requires that the structure of the conclusion be
reviewable against a calibrated reference, not just internally
consistent.

\lean{LOGICAL} names the calibrated structure that lets claims be
compared against each other and against a reference. In the current
code, \lean{LOGICAL} carries \lean{ekg : Calibration.EKG}: the
calibration certificate is part of the class. Logic is not
free-floating cleverness. It carries a scale. The reasoning may be
intricate; the scale tells the reader what the reasoning has been
calibrated against.

The flight's decision rules are logical because each rule names what
measurement triggers what action against what threshold: if fuel below
reserve, divert; if weather below minimums, hold or divert; if traffic
conflict at clearance level, request a different level. Each rule is
calibrated against a published standard and a current measurement.

A chess clock rule is logical because it names when time loss follows
from a measured clock state, not from the players' persuasion. The
clock has the scale; the rule applies it.

An argument that never stops is not yet an answer.

\section{\lean{HALTED}}
\label{se:halted}

Stopping matters. A process that always wants more data, more
deliberation, or more confirmation never licenses action. The
instrument has to be able to commit to an output and treat the output
as halted.

\lean{HALTED} does not say the process was easy. It does not say the
process was short. It says the process reached a form that can be
treated as output. The output may be revised by a later,
separately-named process. But the current process is closed. The
closure is recorded. The next stage can act on the closed output
without re-running the closed process.

The flight halts at landing and the post-flight report is filed at the
gate. Further analysis can continue, but this flight has an actionable
result: the airframe is released for the next leg or held for
maintenance, the crew is rotated or rested, the flight data is
downloaded. The halt is public and dated.

A bake timer halts the process enough to check the loaf. The bread may
rest afterward, but the timer creates an output moment the cook can
act on.

A halted answer must cross the instrument.

\section{\lean{MEASURED}}
\label{se:measured}

An answer must be taken by an instrument. Confidence is not
measurement. A halted output that has not crossed a measurement
process is opinion delivered with timing.

\lean{MEASURED} says the halted output has crossed the measurement
process under the named instrument's discipline. The crossing produces
a typed record: not just ``the answer'' but ``the answer as the
instrument read it.'' The next reader can challenge the instrument or
the measurement, but the instrument's pass through the output is a
public event the reader cannot ignore.

The measured answer is the written flight report with altitudes flown,
fuel burn, deviations, ATC interactions, mechanical notes, and pilot
comments. The controller can act because the answer crossed the named
instruments: the flight data recorder, the fuel system, the ATC log,
and the maintenance database. A confident pilot's verbal assurance is
not a substitute for the measured record.

A piano tuner does not merely believe the note is close. The pitch
crosses an ear, a reference, and often a meter. The tuner's belief is
useful; the crossing is what gets logged.

A measured structure must survive translation.

\section{\lean{COMPILED}}
\label{se:compiled}

Form must survive machinery. The measured record was produced in one
instrument's idiom. The next stage uses a different idiom. The
translation is itself a discipline.

\lean{COMPILED} says the structure has survived translation into the
form the next stage can actually execute. The translation is not a
copy. It is a re-rendering that preserves the structural commitments
while changing the surface. A compilation that loses commitments is
not a compilation; it is a different document. The record retains its
provenance through the translation.

The flight log, the ATC log, the maintenance entries, the fuel slips,
the crew duty record, and the passenger manifest compile into a single
post-flight record the airline can actually use. The record is queried
for billing, for scheduling, for safety analysis, and for regulatory
filing without any of those queries re-running the flight.

A manuscript becomes a typeset book when chapters, footnotes, figures,
page numbers, and index survive translation into a printed form. The
content has changed machinery without losing structure.

Once compiled, the record permits inference.

\section{\lean{INFERRED}}
\label{se:inferred}

What may now be concluded? The procession has crossed thirty-six
gates. Each gate added a discipline. Each gate refused something the
public might have been tempted to accept without it. The final gate
is the one the procession was built to license: a public sentence
that the record now permits and that the speaker can defend without
claiming more than the record carries.

\lean{INFERRED} is the quiet final name. It does not mean guessed. It
does not mean felt. It does not mean authorized by status. It does
not mean \lean{TRUTH}, which has already appeared in the procession
and has done its work as one obligation among the others.
\lean{INFERRED} means the record now permits the conclusion. The
conclusion is bounded by the record's scope; it does not own the
future, the broader population, or any conclusion the record was not
built to support. A speaker who infers from the record can be
challenged at any of the thirty-six prior gates, and the challenge
will be specific: which witness was missing, which scale was wrong,
which gate let an inadmissible mark through. The inferred sentence is
honest because the gates that licensed it are inspectable. It is also
limited, and the limitation is part of the sentence.

The controller may infer ``flight cleared, flown, and closed under
the record on file.'' The buyer may infer ``buy at this price,'' ``buy
only after repair credit of \$1{,}500,'' or ``walk away.'' The
inferred sentence is practical; the speaker can act on it; and the
action is licensed by what the procession has admitted, measured,
witnessed, localized, halted, and compiled. Neither sentence promises
the future. The aircraft may go to maintenance tomorrow. The car may
need a transmission next year. The inference is current; tomorrow's
record will ask the procession again, and the procession will run
again, and the sentence that comes out then will be the sentence that
the new record permits.

A jury verdict is an inference from admitted evidence under
instructions, not a possession of all truth about what happened. The
verdict licenses action. The action is bounded by the trial's scope.
A different trial, under different evidence and instructions, could
license a different verdict; that does not make either verdict false.

The procession ends here. The next chapter asks what remains when
all the necessary names have been spoken.
